\newpage{\ }
\thispagestyle{empty}

\begin{spacing}{1.0} % Espaciado 1.0 solo aquí

\chapter*{Resumen}
\addcontentsline{toc}{chapter}{Resumen}

\vspace{.5cm}

El norte de la Patagonia Andina presenta un paisaje sujeto a fuego. Con el cambio climático se espera que en esta región aumente la temperatura y disminuya la precipitación, lo cual aumentaría notablemente la incidencia de incendios. En este contexto es de vital importancia avanzar en el entendimiento sobre los diversos factores que controlan el fuego, para así poder anticiparnos a los cambios que se esperan en el régimen de incendios. Para ello resulta necesario, además, desarrollar modelos que representen y simulen estos procesos para así proyectar el régimen de incendios ante condiciones ambientales cambiantes. En esta tesis estudiamos cómo los factores físicos (clima y topografía), la vegetación y la presencia humana controlan el fuego en esta región, desarrollamos modelos que permiten simular el régimen de incendios, y realizamos proyecciones de incidencia del fuego bajo escenarios de clima futuro. Previo al desarrollo de modelos de simulación, abordamos el estudio de los controles físicos y bióticos del fuego a dos escalas. A escala de sitio, evaluamos cómo el contrastante microclima de comunidades vegetales aledañas regula la humedad del combustible, contribuyendo a diferencias en inflamabilidad. Y a escala de paisaje, analizamos cómo el clima y la topografía regulan la incidencia del fuego interactuando con la vegetación.

En el norte de la Patagonia Andina dos tipos de comunidades dominan el paisaje: matorrales, de relativamente alta inflamabilidad y dominados por especies rebrotantes, y bosques, de menor inflamabilidad y con baja resiliencia al fuego. Sus diferencias en inflamabilidad se deben en gran parte a su contrastante estructura: los matorrales presentan mayor cantidad y continuidad de combustibles finos. Comparando bosques y matorrales aledaños, encontramos que, además, los bosques presentan un microclima más frío y húmedo que los matorrales, en donde la humedad de combustibles muertos es mayor. De esta manera, su menor inflamabilidad se puede explicar por su menor cantidad y continuidad de combustibles finos, pero también por una mayor humedad del combustible muerto. 

La menor inflamabilidad de los bosques también se refleja en una menor incidencia del fuego cuando el proceso se estudia a escala de paisaje. Sin embargo, a esta escala, otros factores pueden confundir el efecto de la vegetación. Por ejemplo, los bosques tienden a ocupar las zonas de mayor precipitación y altitud (menor temperatura), y los matorrales ocupan las zonas más cercanas a poblados, expuestos a una mayor frecuencia de igniciones. Para poder modelar el régimen de incendios controlando por las variables adecuadas, evaluamos cómo la topografía y la precipitación media anual controlan el fuego interactuando con la vegetación. Mapeamos todos los incendios ocurridos en un período de 24 años (1999-2022) en el norte de la Patagonia Andina (39° a 44.5° S), y evaluamos cómo variaba la probabilidad de que las celdas se hubieran quemado en todo el período en función de características ambientales. Encontramos que la incidencia de incendios disminuye con la altitud y con la precipitación, pero el efecto de la precipitación en gran parte se debe a cómo se ordenan comunidades de estructura contrastante a lo largo del gradiente: los tipos de vegetación con mayor cantidad y continuidad de combustible fino y menor humedad intrínseca del combustible se presentaron en zonas de precipitación baja e intermedia. A su vez, detectamos un claro efecto unimodal de la productividad, tanto dentro como entre comunidades. 

Partiendo del entendimiento sobre cómo distintos factores controlan la incidencia del fuego desarrollamos un simulador del régimen de incendios, compuestos por tres modelos espacialmente explícitos: ignición, escape y propagación. Los modelos fueron definidos en tiempo y espacio discretos, con una resolución espacial de 30 m y temporal de 14 días. Las igniciones pueden ser causadas por rayos o por actividad humana, y una vez que ocurren pueden quemar completamente la celda original y escapar de ella. Dado el escape, el modelo de propagación regula cómo se contagia fuego celda a celda. En todos los modelos se incluyeron efectos de las condiciones atmosféricas, utilizando el índice de peligrosidad Fire Weather Index, la vegetación (tipo de vegetación y su productividad) y la topografía. Además, en los modelos de ignición por humanos y escape se consideró la distancia al poblado y al camino más cercanos, y en el modelo de propagación se incluyó el efecto del viento. Para estimar los modelos de ignición y escape utilizamos los registros de focos del Parque Nacional Nahuel Huapi, mientras que el modelo de propagación fue estimado utilizando los polígonos de incendios mapeados para el objetivo anterior. En los tres modelos hallamos un claro efecto del clima, indicando que a mayor Fire Weather Index (mayor temperatura, menor humedad del aire, menor precipitación y mayor velocidad del viento) la probabilidad de ignición, escape y propagación aumentan. Encontramos que, si bien las igniciones humanas son más frecuentes que las causadas por rayos, las últimas tienen mayor probabilidad de escape por ocurrir lejos de caminos, en donde el control temprano es dificultoso. Con estos modelos proyectamos la incidencia del fuego en el Parque Nacional Nahuel Huapi para los períodos 2040-2049 y 2090-2099, considerando cuatro escenarios climáticos contrastantes. Las simulaciones indican que para 2040-2049 la probabilidad de quema anual aumentará en un factor de entre 2.07 y 2.61, mientras que para 2090-2099 se esperan aumentos en factores de entre 1.36 y 31.17, dependiendo del escenario climático.

Nuestro trabajo señala la importancia de implementar políticas eficientes de manejo del fuego para atenuar sus efectos sobre las sociedades y los ecosistemas. En particular, con el aumento de la incidencia del fuego se espera que los bosques andino-patagónicos se retraigan, a expensas de un aumento en la cobertura de matorrales y pastizales. Estas transiciones, a su vez, promoverían aún más la ocurrencia de incendios. 

\noindent \textbf{Palabras clave:}  \textit{Probabilidad de quema, modelos, propagación, vegetación, topografía.}


\end{spacing}